\documentclass[a4paper,12pt]{article}
\usepackage[T2A]{fontenc}
\usepackage[utf8]{inputenc}
\usepackage[russian]{babel}
\usepackage{amsmath,amssymb,mathtools}
\usepackage{helvet}
\renewcommand{\familydefault}{\sfdefault}
\usepackage{icomma}
\usepackage[a4paper,margin=20mm,headheight=25pt,headsep=9pt,footskip=12mm]{geometry}
\usepackage{microtype}
\usepackage{tikz}
\usetikzlibrary{calc,arrows.meta,positioning,shapes.geometric}
\usepackage{booktabs,tabularx,array}
\usepackage{enumitem}
\usepackage{xcolor}
\usepackage{fancyhdr}
\usepackage{setspace}
\usepackage{hyperref}
\usepackage{parskip}
\definecolor{accent}{HTML}{087C80}
\definecolor{accentdark}{HTML}{14565B}
\definecolor{ink}{HTML}{28363C}
\definecolor{subtle}{HTML}{6A7A7C}
\definecolor{redink}{HTML}{A84637}
\hypersetup{colorlinks=true,linkcolor=accentdark,urlcolor=accent}
\setstretch{1.15}
\setlength{\parskip}{6pt}
\setlist[enumerate]{leftmargin=*,itemsep=9pt,topsep=4pt}
\setlist[itemize]{leftmargin=*,itemsep=3pt,topsep=3pt}
\pagestyle{fancy}
\fancyhf{}
\fancyhead[L]{\small\textcolor{subtle}{ОПТИМИЗАЦИЯ ФУНКЦИЙ}}
\fancyhead[R]{\small\textcolor{subtle}{Николь Саркисьянц}}
\fancyfoot[C]{\small\textcolor{subtle}{\thepage}}
\renewcommand{\headrulewidth}{0.35pt}
\renewcommand{\headrule}{\hbox to\headwidth{\color{accent!30}\leaders\hrule height \headrulewidth\hfill}}
\newcommand{\idea}[1]{\par\noindent\textcolor{accentdark}{\textbf{Главная идея.}}\ #1\par}
\newcommand{\note}[1]{\par\noindent\textcolor{redink}{\textbf{Проверьте себя.}}\ #1\par}
\newcommand{\ex}[1]{\par\noindent\textcolor{accentdark}{\textbf{Пример.}}\ #1\par}
\newcommand{\task}[2]{\item \textbf{#1}\\[-1pt]#2}
\begin{document}
\thispagestyle{empty}
\begin{center}
\vspace*{19mm}
{\small\bfseries\textcolor{accent}{МАТЕМАТИКА}}\par
\vspace{27mm}
{\fontsize{34}{39}\selectfont\bfseries Чек-лист}\par
\vspace{12mm}
{\LARGE\bfseries Оптимизация функций}\par
\vspace{4mm}
{\Large Производная, экспонента и работа с константами}\par
\vspace{20mm}
{\normalsize Подготовка к ЕГЭ по профильной математике}\par
\vspace{21mm}
{\Large\bfseries Николь Саркисьянц}\par
\vspace{7mm}
29 октября 2026 г.\par
\vfill
{\small Лёвин Артём Александрович}\par
{\small эксклюзивно для Николь Саркисьянц}\par
\vspace{12mm}
\begin{tikzpicture}
\draw[accent!60,line width=1pt] (-6.6,0) .. controls (-4.5,0.7) and (-2.3,-0.6) .. (0,0)
  .. controls (2.4,0.6) and (4.5,-0.4) .. (6.6,0.06);
\end{tikzpicture}
\end{center}
\clearpage

\section*{1. Что мы ищем в задачах на оптимизацию}
\idea{Если требуется найти, при каком значении величины функция достигает наибольшего или наименьшего значения, сначала определите переменную и её допустимые значения.}

\begin{itemize}
\item Запись $f(v)$ подсказывает, что изменяется $v$; остальные буквы могут быть заданными константами.
\item Постоянные коэффициенты удобно объединять: сложная формула после замены часто превращается в короткое произведение.
\item Производная показывает, на каких промежутках функция возрастает и убывает; найденные стационарные точки необходимо проверить по знакам производной и границам области.
\end{itemize}

\ex{Пусть при фиксированных $M,R,T>0$ дана функция
\[
 f(v)=4\pi v^2\left(\frac{M}{2\pi RT}\right)^{3/2}
   e^{-\frac{Mv^2}{2RT}},\qquad v\ge0.
\]
Здесь переменная --- $v$. Обозначим
\[
 A=4\pi\left(\frac{M}{2\pi RT}\right)^{3/2}>0,
 \qquad B=-\frac{M}{2RT}<0.
\]
Тогда $f(v)=Av^2e^{Bv^2}$. Именно такую форму удобнее дифференцировать.}

\section*{2. Четыре правила, которые нужны чаще всего}
\begin{tabularx}{\linewidth}{@{}p{0.36\linewidth}X@{}}
\toprule
\textbf{Правило} & \textbf{Запись} \\ \midrule
Произведение & $(uv)'=u'v+uv'$ \\[3pt]
Сложная экспонента & $(e^{g(x)})'=g'(x)e^{g(x)}$ \\[3pt]
Корень & $(\sqrt{x})'=\dfrac{1}{2\sqrt{x}}$ при $x>0$ \\[5pt]
Частное & $\displaystyle\left(\frac{u}{v}\right)'=\frac{u'v-uv'}{v^2}$ при $v\ne0$ \\ 
\bottomrule
\end{tabularx}

\ex{Для $g(v)=Bv^2$, где $B$ постоянно, имеем $g'(v)=2Bv$, поэтому
\[
\boxed{\bigl(e^{Bv^2}\bigr)'=2Bv\,e^{Bv^2}.}
\]
\textbf{Обязательно} сохраняйте множитель $e^{Bv^2}$ после дифференцирования показателя.}

\section*{3. Вынесение общего множителя}
Можно вынести множитель за скобки, только если \textbf{каждое слагаемое} содержит этот множитель.
\[
2x^2+3x+5=\underbrace{x(2x+3)}_{\text{первые два слагаемых}}+5.
\]
\note{Равенство $2x^2+3x+5=x(2x+3+5)$ неверно: у числа $5$ нет множителя $x$.}

\section*{4. Разобранный тип I: энергия молекулы}
В задаче с плотностью распределения кинетической энергии встречается функция
\[
 f(\varepsilon)=\frac{2}{\sqrt\pi}\left(\frac1{kT}\right)^{3/2}
 \sqrt\varepsilon\,e^{-\varepsilon/(kT)},\qquad \varepsilon\ge0,
\]
где $k,T>0$ --- постоянные. Обозначив
\[
 A=\frac2{\sqrt\pi}\left(\frac1{kT}\right)^{3/2}>0,
 \quad B=\frac1{kT}>0,
\]
получим $f(\varepsilon)=A\sqrt\varepsilon\,e^{-B\varepsilon}$.

\textbf{Шаг 1.} При $\varepsilon>0$ применяем правило произведения:
\begin{align*}
 f'(\varepsilon)
 &=A\frac1{2\sqrt\varepsilon}e^{-B\varepsilon}
 -AB\sqrt\varepsilon\,e^{-B\varepsilon}\\
 &=A e^{-B\varepsilon}\left(\frac1{2\sqrt\varepsilon}-B\sqrt\varepsilon\right)\\
 &=\frac{Ae^{-B\varepsilon}}{2\sqrt\varepsilon}\,(1-2B\varepsilon).
\end{align*}

\textbf{Шаг 2.} Первые множители положительны при $\varepsilon>0$, поэтому знак производной определяется выражением $1-2B\varepsilon$. Приравниваем его к нулю:
\[
1-2B\varepsilon=0
\quad\Longrightarrow\quad
\varepsilon_*=\frac1{2B}=\boxed{\frac{kT}{2}}.
\]

\textbf{Шаг 3.} На числовой прямой отмечаем $0$ и $\varepsilon_*$.
\begin{center}
\begin{tikzpicture}[>=Latex,scale=1]
\draw[->,accent,line width=.9pt] (0,0)--(12,0) node[right] {$\varepsilon$};
\draw (1,-.12)--(1,.12) node[below=5pt] {$0$};
\draw (6,-.12)--(6,.12) node[below=5pt] {$\varepsilon_*=\dfrac{kT}{2}$};
\node[accentdark] at (3.2,.42) {$f'(\varepsilon)>0$};
\node[redink] at (9,.42) {$f'(\varepsilon)<0$};
\draw[->,line width=1.1pt,accent] (1.7,-.65)--(5.7,-.14);
\draw[->,line width=1.1pt,redink] (6.3,-.14)--(10.3,-.65);
\end{tikzpicture}
\end{center}
Слева функция возрастает, справа убывает. При $\varepsilon=0$ значение равно нулю; при $\varepsilon\to+\infty$ оно стремится к нулю. Значит, в $\varepsilon_*$ достигается \textbf{наибольшее} значение.

\ex{Если $k=1{,}38\cdot10^{-23}$ Дж/К и $T=80$ К, то
\[
\varepsilon_*=\frac{1{,}38\cdot10^{-23}\cdot80}{2}
=\boxed{5{,}52\cdot10^{-22}\text{ Дж}}.
\]
Сначала сокращаем символические выражения, затем подставляем числа.}
\section*{5. Разобранный тип II: скорость молекулы}
После объединения констант получаем функцию
\[
 f(v)=Av^2e^{Bv^2},\qquad A>0,\ B<0,\ v\ge0.
\]
Разложим её на множители $u=Av^2$ и $w=e^{Bv^2}$:
\[
 u'=2Av,\qquad w'=2Bv\,e^{Bv^2}.
\]
Тогда
\begin{align*}
 f'(v)&=2Av\,e^{Bv^2}+Av^2\cdot2Bv\,e^{Bv^2}\\
 &=2Av\,e^{Bv^2}\underbrace{(1+Bv^2)}_{\text{определяет знак при }v>0}.
\end{align*}

Решаем $f'(v)=0$. Первый вариант --- $v=0$ --- граница области. На $v>0$ требуется
\[
1+Bv^2=0\quad\Longrightarrow\quad
 v^2=-\frac1B\quad\Longrightarrow\quad
 \boxed{v_*=\sqrt{-\frac1B}}.
\]
Знак $1+Bv^2$ положителен до $v_*$ и отрицателен после $v_*$. Поэтому найденная положительная скорость даёт наибольшее значение функции.

Возвращаем постоянные $B=-M/(2RT)$:
\[
\boxed{v_*=\sqrt{\frac{2RT}{M}}}\qquad
(R,T,M>0).
\]
\note{В скобках после дифференцирования должно быть $1+Bv^2$. Следовательно, решается уравнение для $v^2$, а не для $v^3$. Величина $B$ уже содержит знак минус.}

\section*{6. Когда подстановка констант не помогает}
Рассмотрим характерную форму, встречавшуюся в задачах о колебаниях:
\[
F(\omega)=\frac{UL\omega}
{\sqrt{R^2+\left(L\omega-\frac1{C\omega}\right)^2}},\qquad\omega>0.
\]
Здесь $U,L,R,C$ постоянны, однако $\omega$ встречается и в числителе, и под корнем, и в знаменателе внутренней дроби. Обозначения $A,B$ почти не упростят производную. Полезный маршрут:
\begin{enumerate}
 \item Укажите ОДЗ: $\omega>0$ (в частности, $\omega\ne0$).
 \item Найдите производную выражения под корнем по правилу сложной функции.
 \item Примените правило частного или перепишите деление как умножение на степень $-\frac12$.
 \item Вынесите общий множитель только там, где он действительно присутствует.
\end{enumerate}
\idea{Начинайте вычисления с буквами. Числа подставляйте тогда, когда выражение уже упрощено.}

\textbf{Как выглядит прямое дифференцирование.}
Пусть $h(\omega)=L\omega-\frac1{C\omega}$,
$Q(\omega)=R^2+h(\omega)^2$; тогда
\[
 h'(\omega)=L+\frac1{C\omega^2},\qquad
 Q'(\omega)=2h(\omega)h'(\omega).
\]
Перепишем $F(\omega)=UL\omega\,Q(\omega)^{-1/2}$. Правило
произведения вместе с правилом сложной функции даёт
\begin{align*}
 F'(\omega)
 &=UL\left(\frac1{\sqrt Q}-\frac{\omega Q'}{2Q^{3/2}}\right)\\
 &=\frac{UL\bigl(Q-\omega h h'\bigr)}{Q^{3/2}}.
\end{align*}
Здесь $Q=Q(\omega)$, $h=h(\omega)$. Только теперь раскрываем
скобки и сокращаем подобные члены:
\[
\boxed{F'(\omega)=\frac{UL\left(R^2-\dfrac{2L}{C}
   +\dfrac{2}{C^2\omega^2}\right)}
 {\left(R^2+\left(L\omega-\dfrac1{C\omega}\right)^2\right)^{3/2}}.}
\]
Это пример, когда производная частного и сложной функции
полезнее случайных числовых подстановок.
\clearpage

\thispagestyle{fancy}
\begin{center}
{\LARGE\bfseries Алгоритм поиска максимума}\par
\vspace{4mm}
{\small Для функций с произведением степенного и показательного множителей}\par
\end{center}
\vspace{6mm}
\begin{center}
\begin{tikzpicture}[
 node distance=10mm,
 start/.style={ellipse,draw=accent,line width=.9pt,align=center,minimum width=56mm,minimum height=10mm,fill=white},
 stage/.style={rectangle,rounded corners=2pt,draw=accent!70,line width=.8pt,align=center,text width=93mm,inner sep=8pt,fill=white},
 decision/.style={diamond,aspect=2.5,draw=accentdark,line width=.9pt,align=center,inner sep=3pt,text width=41mm,fill=white},
 arr/.style={-{Latex[length=3mm]},line width=.9pt,draw=accentdark}
]
\node[start] (start) {Начало};
\node[stage,below=of start] (var) {Выделить переменную и ОДЗ; отделить постоянные коэффициенты};
\node[decision,below=11mm of var] (choice) {Замена упрощает формулу?};
\node[stage,below left=17mm and 7mm of choice,text width=35mm] (yes) {Ввести постоянные $A,B$};
\node[stage,below right=17mm and 7mm of choice,text width=35mm] (no) {Оставить исходные коэффициенты};
\node[stage,below=37mm of choice] (der) {Найти $f'(x)$: произведение, цепное правило, частное};
\node[stage,below=of der] (zeros) {Решить $f'(x)=0$; проверить допустимость найденных точек};
\node[stage,below=of zeros] (signs) {Исследовать знак $f'$ слева и справа; учесть границы ОДЗ};
\node[start,below=of signs] (end) {Записать максимум и ответ};
\draw[arr] (start)--(var);
\draw[arr] (var)--(choice);
\draw[arr] (choice) -| node[pos=.18,above left]{\small да} (yes);
\draw[arr] (choice) -| node[pos=.18,above right]{\small нет} (no);
\draw[arr] (yes.south) |- ($(der.north)+(0,5mm)$) -- (der.north);
\draw[arr] (no.south) |- ($(der.north)+(0,5mm)$);
\draw[arr] (der)--(zeros);
\draw[arr] (zeros)--(signs);
\draw[arr] (signs)--(end);
\end{tikzpicture}
\end{center}
\clearpage

\section*{7. Типичные ошибки и самопроверка}
\begin{enumerate}
 \item \textbf{Потеря экспоненты.} Правильно: $(e^{Bv^2})'=2Bv\,e^{Bv^2}$.
 \item \textbf{Путаница степеней при вынесении.} Из $2Av+2ABv^3$ выносим $2Av$ и получаем $2Av(1+Bv^2)$.
 \item \textbf{Общий множитель отсутствует у одного слагаемого.} Тогда выносить его из всей суммы нельзя.
 \item \textbf{Неверный знак константы.} Для $e^{-Mv^2/(2RT)}$ удобно брать $B=-M/(2RT)<0$.
 \item \textbf{Найден только корень производной.} Обязательно убедитесь, что знак $f'$ меняется с плюса на минус; проверьте границы области.
 \item \textbf{Слишком ранняя подстановка чисел.} Сохраняйте буквенный вид, чтобы не потерять степень десяти и не усложнять дроби.
\end{enumerate}
\idea{Итоговая проверка: названа переменная; указана ОДЗ; производная посчитана по правилам; множители раскрыты и вынесены корректно; доказан максимум; ответ имеет нужные единицы измерения.}

\clearpage
\section*{Тренировочный блок --- Путь А}
\setlength{\abovedisplayskip}{5pt}\setlength{\belowdisplayskip}{5pt}
\textcolor{subtle}{10 заданий, от отработки правил к задачам на оптимизацию. Решения запишите последовательно; в задачах на максимум обоснуйте изменение знака производной.}
\begin{enumerate}[itemsep=3pt,topsep=3pt]
\task{Определите переменную и коэффициенты.}{Дана функция
\[
F(v)=3v^2\left(\frac{m}{4\pi RT}\right)^{3/2}e^{-\frac{mv^2}{4RT}},\quad v\ge0,
\]
где $m,R,T>0$ фиксированы. Запишите её в виде $F(v)=Av^2e^{Bv^2}$. Найдите $A$, $B$ и укажите знак каждого.}
\task{Вынесите общий множитель.}{Представьте выражение $4x^2+6x+7$ в виде $x(\ldots)+7$. Объясните, почему нельзя вынести $x$ из всей суммы.}
\task{Производная сложной функции.}{Вычислите производную $y=e^{-3x^2}$.}
\task{Корень и экспонента.}{Найдите производную $y=4\sqrt{x}\,e^{-x/3}$ при $x>0$ и вынесите общий показательный множитель.}
\task{Два правила одновременно.}{Найдите и разложите на множители производную $y=5v^2e^{-v^2/8}$.}
\task{Простейшая оптимизация.}{Найдите точку наибольшего значения функции $y=xe^{-2x}$ при $x\ge0$.}
\task{Максимум с корнем.}{Найдите точку наибольшего значения $y=\sqrt{x}\,e^{-x/6}$ при $x\ge0$.}
\task{Производная с корнем в знаменателе.}{Вычислите и упростите производную функции $y=\dfrac{t}{\sqrt{t^2+4}}$.}
\task{Энергия молекулы.}{Функция плотности имеет вид $f(\varepsilon)=A\sqrt\varepsilon\,e^{-\varepsilon/(kT)}$ при $A,k,T>0$ и $\varepsilon\ge0$. Найдите энергию максимума при $k=1{,}38\cdot10^{-23}$ Дж/К, $T=80$ К.}
\task{Скорость молекулы.}{При $v\ge0$ функция задана формулой
\[
 f(v)=4\pi v^2\left(\frac{M}{2\pi RT}\right)^{3/2}e^{-Mv^2/(2RT)}.
\]
Найдите скорость, при которой она максимальна, если $R=8{,}31$ Дж/(моль$\cdot$К), $T=450$ К, $M=0{,}004$ кг/моль. Округлите до целого м/с.}
\end{enumerate}
\clearpage
\section*{Ответы к тренировочному блоку}
\renewcommand{\arraystretch}{1.42}
\begin{tabularx}{\linewidth}{@{}p{.075\linewidth}X@{}}
\toprule
\textbf{№} & \textbf{Ответ} \\ \midrule
1 & $A=3\left(\dfrac{m}{4\pi RT}\right)^{3/2}>0$, $B=-\dfrac{m}{4RT}<0$; переменная $v$.\\
2 & $x(4x+6)+7$; число $7$ не содержит множителя $x$.\\
3 & $y'=-6xe^{-3x^2}$.\\
4 & $y'=4e^{-x/3}\left(\dfrac1{2\sqrt{x}}-\dfrac{\sqrt{x}}3\right)$.\\
5 & $y'=10ve^{-v^2/8}\left(1-\dfrac{v^2}{8}\right)$.\\
6 & $x_*=\dfrac12$.\\
7 & $x_*=3$.\\
8 & $y'=\dfrac{4}{(t^2+4)^{3/2}}$.\\
9 & $\varepsilon_*=\dfrac{kT}{2}=5{,}52\cdot10^{-22}$ Дж.\\
10 & $v_*=\sqrt{\dfrac{2RT}{M}}\approx1367$ м/с.\\
\bottomrule
\end{tabularx}
\vspace{10mm}
\textcolor{accentdark}{\textbf{Последняя самопроверка}}\par
Если в условии спрашивают \emph{наибольшее значение функции}, после нахождения точки максимума нужно дополнительно подставить эту точку в исходную функцию. Если спрашивают, \emph{при каком значении переменной} достигается максимум, достаточно найденной точки с обоснованием.
\end{document}
